\fancysection{Objectives and Significance}
%
\label{sec:objectives_and_significance}

We propose a focused project to combine \tess{} photometry from both high
cadence stamps as well as full frame images (FFIs), with spectral energy
distribution (SED) information to provide a catalog of detached eclipsing
binaries that will:

\begin{enumerate}
%
    \item Fully characterize the often complex and highly correlated uncertainty
        distributions of their physical parameters allowed by the available
        observations
%
    \item Measure spin and differential rotation of \tess{} EB primary stars
%
    \item Detect ellipsoidal variations, reflected light, and Doppler beaming
        effects
%
    \item Measure eclipsing timing variability (ETV) (e.g. produced by apsidal
        precession or extra objects in the system).
%
\end{enumerate}

The proposed catalog will provide the required input for numerous investigations
in stellar and exoplanetary physics.

\subsection{Expected Applications}
%
\label{sec:applications}

The proposed detailed catalog of thousands of EBs will serve as the bedrock of
projects our group and others will pursue after this proposal to measure the
rate at which tidal energy is converted to heat and the associated torques in
short period binary stars and exoplanet systems. The features of the proposed
catalog provide unique advantages to such studies. For example, the catalog
will include populations of systems of comparable size for which tides will:
strongly affect the orbit and the spin of the components, leave the orbit mostly
unaffected but strongly affect the spins, are weak enough to have no observable
signatures. These populations provide natural controls to each other when
investigating tidal dissipation allowing the validation of the modeling of tidal
and non-tidal effects for very similar objects, or even isolating multiple
effects within a single object (e.g. tidal effects on orbital eccentricity
    generally probes dissipation of tidal waves of higher frequency than the
tidal effects on the spin). Stellar tides, in addition to directly being a
crucial factor shaping the observed exoplanet population, are also required
input to measuring tides in planets, since the latter are always accompanied by
the former. Since the proposed catalog will probe similar population of stars to
exoplanet hosts, the measured stellar tides will be directly applicable when
modeling exoplanet systems.

In addition to tides, binary stars, especially in combination with spin
measurements, provide a unique laboratory for testing a host of additional
stellar physics, like providing direct tests of stellar evolution and
asteroseismology \citep[e.g.][]{delBurgo_AllendePrieto_18, Gaulme_et_al_16},
magnetic breaking \citep[e.g.][]{ElBadry_et_al_22}, limb darkening and stellar
atmospheres modeling \citep[e.g.][]{<++>}, stellar differential rotation and
magnetic activity \citep[e.g.][]{Reinhold_et_al_13, Jermyn_et_al_20a,
Jermyn_et_al_20b}. In addition to observational signatures probing the surface
of the star, measurements of apsidal precession can constrain the tidal love
number and hence provide one of very few direct probes of internal stellar
structure.

Apart from stellar physics, EBs are the most common source of false positive
signals in transiting exoplanet searches and understanding their properties and
prevalence is crucial to interpreting exoplanet observations
\citep[e.g.][]{}<++>. Eclipsing binaries are also important inputs to fields that
at first glance appear unrelated, like Spectral synthesis for galaxy science
\citep[e.g.][]{Eldridge_et_al_17} or even studying the galactic dark matter
\citep[e.g.][]{Chakrabarti_22}.

\subsection{Comparison to Existing EB Catalogs}
%
\label{sec:comparison_to_existing}
%
While the proposed catalog is not the first collection of EBs to be produced out
of the recent flood of transiting planet surveys, including \tess{}, it has
unique characteristics to all publications to date that make it particularly
valuable for the science cases outlined above.

Broadly speaking, catalogs using ground based photometry do not come anywhere
close to the level of precision of \tess{} or \kepler{}, which in turn
dramatically limits the attainable precision of the inferred physical
parameters. For example, both \tess{} and \kepler{} lightcurves allow measuring
the time of primary and secondary eclipse with sufficient precision to determine
$e\cos\omega$ (orbital eccentricity times cosine of the longitude of periapsis)
to a precision of $10^{-4}$ and often much better, where as typical
uncertainties from ground based data are orders of magnitude larger
\citep[c.f.][]{}<++>. Space-based photometry offers similar boost in precision
in the determination of most of the physical parameters of the systems.
Furthemore, effects often not detectable from the ground are easily detectable
from space (e.g. ellipsoidal variations, Doppler boosting etc.) which in turn
provide crucial new constraints on the binary parameters and enable new science
use cases.

On the other hand, this effort will produce the largest and best characterized
collection of EBs from space photometry to date. Thanks to the full sky coverage
of \tess{}, those will be the brightest EBs available, making followup
observations of individual systems as easy as possible should any be necessary
for a particular application. We compare our expected results to the state of
the art \tess{} and \kepler{} catalogs available in the literature below.

\citep{Prsa_et_al_22} and \citep{Justesen_Albrecht_21} both used \tess{}
photometry to search for eclipsing binary stars and \citep{Justesen_Albrecht_21}
combined photometry with SED information to find maximum likelihood parameters
for each binary. Our proposal would improve on both efforts in four crucial
ways.  First, both studies made no attempt to provide error bars for the
extracted system parameters. We propose to apply Bayesian analysis (BA) to fully
characterize the uncertainty distributions of the parameters allowed by the
observational data.  This is a crucial requirement for many of the applications
described above so that for example uncertainties in the inferred tidal
dissipation can be characterized, or any suggested tension between stellar
structure and evolution or stellar atmosphere models can be quantified. Second,
\citep{Justesen_Albrecht_21} only used data from the first 13 sectors from
\tess{}, while \citep{Prsa_et_al_22} used 26. In contrast, up to March 20, 2023
\tess{} had completed observations of 62 sectors, revisiting the same locations
on the sky analyzed by \citep{Prsa_et_al_22} and \citep{Justesen_Albrecht_21} a
second or even a third time in some cases as well as adding new locations. The
extended time coverage will dramatically boost the sensitivity of this effort to
longer orbital or spin periods, significantly improve the precision of the
parameter determination, provide important opportunities to filter out false
positives from automatically selected EB or rotational variability candidates,
and provide much larger leverage for detecting eclipse timing variations.
Third, both prior studies relied only on \tess{} 2-min cadence stamp data. While
we also plan to begin by focusing on stamp data, we will also fully exploit the
FFIs, increasing the number of lightcurves search by a factor of $\sim50$, and
hence by a comparable (though likely somewhat smaller) factor the number of
expected detections. Finally, none of the existing \tess{} EB collections come
with spin measurements. While efforts to measure spins using \tess{} photometry
exist \citep[e.g.]{}<++>, none focused on eclipsing binaries, which require
specialized analysis to disentangle the rotational variability from the eclipses
and other binary driven variability. Spin measurements are very valuable for the
tidal studies described above, enabling the cross-checks and validations
described, and are required for many of the other applications (studies of
magnetic breaking, differential rotation, magnetic activity etc.).

A combination of two studies of \kepler{} EBs comes closest to the catalog we
plan to create: \citep{Windemuth_et_al_19} who combined photometry and SED and
used Bayesian analysis to fully characterize the physical parameters of 728 EBs
and their uncertainty distributions, and \citep{Lurie_et_al_17} who searched
known \kepler{} EBs lightcurves for signatures of rotational variability,
determining rotational periods for 816 binaries and measuring differential
rotation in many of those. The two catalogs overlap for 363 EBs. Apart from
improved modeling of lightcurves and SED compared to \citep{Windemuth_et_al_19}
and a more rigorous search for signatures of rotational modulation than
\citep{Lurie_et_al_17}, the proposed analysis will likely produce more than
$10^4$ EBs with detailed physical parameters and fully characterized
uncertainties, and for thousands of those, the spin and differential rotation
will be measured (see \refsec{expected_outcomes}). Having an order of magnitude
larger catalog of EBs enables qualitatively different studies of tides. For
example, it is expected that tidal dissipation will be different between stars
with different internal structure (e.g. with and without significant surface
convective zones, or with or without convective core). It is also expected to
change as a function of tidal frequency and amplitude, stellar spin and possibly
other parameters. The larger sample size to come out of the proposed effort is
required to study these dependencies, as smaller datasets quickly run out of
objects when sliced along all these dimensions.

\subsection{Expected Outcomes}
%
\label{sec:expected_outcomes}

\citep{Prsa_et_al_22} searched $\sim2\times10^5$ lightcurves extracted from the
high cadence stamp of the first 26 \tess{} sectors, detecting 4\,584 EBs. Not
all of these binaries will be included in the catalog we plan to produce since
the \citep{Prsa_et_al_22} catalog includes contact binaries and binaries for
which the signal to noise will not be sufficient for high precision physical
parameter determination.  Requiring detached binaries (the sum of the stellar
radii not exceeding half the pericenter distance), primary eclipse depth of at
least 5\%, secondary eclipse depth of at least 1\% allows us to estimate that we
will be able to obtain high quality and high precision characterization of at
least 800 binaries from stamp lightcurves with about double that number
reasonably well constrained. We obtain the same estimate by scaling the results
of \citep{Justesen_Albrecht_21} who found ~350 high precision binaries and more
than 800 with reasonable constraints just from the southern hemisphere. With
double or triple the observations for each binary, we will undoubtedly be able
to extract high precision physical parameters from somewhat shallower eclipses
than \citep{Justesen_Albrecht_21}, and will have improved sensitivity to longer
orbital periods. Therefore, \textbf{to a good approximation our EB catalog based
on stamp data alone should include $\sim1000$ EBs, roughly half of which will
have high precision measurements of their parameters.}

Predicting the expected EB detections from FFIs is less straightforward. We plan
to use the Quick Look Pipeline (QLP) lightcurves based on the full frames. The
QLP provides lightcurves for more than $10^7$ stars brighter than \tess{}
magnitude $T<13.5$. Naively, we will search roughly 50 times more lightcurves
and the 30 min cadence is still plenty to resolve all features of the
lightcurve. AS a result, one could estimate creating a catalog of $5\times10^4$
EBs, half with high precision parameters. However, the $2\times10^5$ stamps
targets were not randomly selected. In fact over 1700 known or suspected EBs
were specifically added to the target list through the TESS Guest Investigator
program. On the other hand, the QLP lightcurve database only includes stars down
to $T<13.5$ whereas stamp targets included much fainter stars. The photometric
precision on a 1-h timescale vs magnitude reported for the stamp-based and
FFI-based lightcurves is almost identical. AS a result, on average the FFI
lightcurves will have higher precision than stamp lightcurves. Regardless, it is
reasonable to expect that \textbf{we will end up with $\sim10^4$ EBs with high
    precision parameters and roughly double that number in total from the FFI
lightcurves.}

In order to get a lower limit to the fraction of binaries for which we also
expect to measure the spin period and often differential rotation, we look to
\citep{CantoMartins_et_al_20}, who searched 1000 \tess{} objects of interest and
measured 131 unambiguous rotation periods (13\%) using less than a year of
observations, and \citep{Oelkers_et_al_18} who searched $\approx4\times10^6$
ground based lightcurves (from the KELT survey) detecting 47\,174 (12\%) with
clear spin signatures with periods under 13 days (a somewhat conservative upper
limit to what can be detected using \tess). The former survey relied on much 2
or 3 times fewer observations per target than we will use and the latter used
ground based data of much lower photometric precision, so we can confidently
conclude 10\%-15\% is a lower limit to the fraction of stars exhibiting
rotational modulations detectable by \tess{} at any given time. When focusing on
EBs that fraction is bound to be higher, since tidal interactions spin-up short
period binaries increasing the fraction with periods under 13 days, and likely
increasing the signal since faster rotators tend to have higher magnetic
activity. An upper limit to the fraction of EBs for which we will measure the
spin, we turn to \citep{Lurie_et_al_17} who used \kepler{} photometry (higher
precision and longer time span than we will have) to measure spin of EBs.
Crossmatching the \citep{Lurie_et_al_17} catalog of firmly detected spins with
the \citep{Windemuth_et_al_19} catalog of EBs with high precision parameters
results in 31\% (227 out of 728 EBs) of the \citep{Windemuth_et_al_19} stars
having a spin period less than 13 days. In short, \textbf{we will be able to
detect spin periods for between 10\% and 30\% of the EBs in our catalog: several
hundred from stamps and many thousands from FFIs.}

In addition to the number of systems, it would be useful to have an idea of the
precision with which various physical parameters will be determined. Due to the
generally highly correlated uncertainties of the parameters, the precision that
is relevant to one application may be very different from another. For example,
tides are sensitive to the ratio of stellar radii to the semimajor axis. This is
usually much better constrained than either of the radii or the semimajor axis
by itself (i.e. marginalized over all other quantities). Regardless,
marginalized uncertainties of individual parameters can still be a useful guide
to the value of the catalog. While \citep{Justesen_Albrecht_21} did not publish
uncertainties, since they were not able to run their analysis long enough to be
fully converged, the order of magnitude of the marginalized uncertainties can be
estimated from the available samples. \textbf{For the high precision EBs
    mentioned above, we expect marginalized uncertainties for the stellar masses
    better than 3\%, and for the radii better than 2\% \citep[private
    communication]{Justesen_Albrecht_21}. For eccentricities, the $e\cos\omega$
    component is much better constrained, with typical uncertainties
$\mathbf{\sim10^{-4}}$, than the $e\sin\omega$ component, with typical
uncertainties $\mathbf{\sim10^{-3}}$ \citep{Justesen_Albrecht_21}.} Measuring
spin periods from spot modulation usually produces extremely precise values.
However, the real uncertainty in that quantity comes from the fact that
different longitudes of the same star rotate at different rates. As a result,
the uncertainty in the location of the stellar spots producing the rotational
modulation dominates the uncertainty with which the rotation period can be
measured. As a result, \textbf{the spin precision will be the best theoretically
possible}, but still limited to $\sim$10\%.
